% -----------------------------------------------------------------------------
% Full scalar-interaction laboratory in strict SI units
% -----------------------------------------------------------------------------

复标量物质场与实标量媒介场提供了不含自旋指标的最小散射模型。在这个体系中，作用量、Dyson 展开、Wick 收缩、LSZ 截肢、散射振幅、相空间和截面可以在同一套归一化下完整闭合。四动量继续采用全书的 SI 约定
\begin{equation*}
 p^\mu=(E_{\mathbf p}/c,\mathbf p),\qquad p^2=m_{\rm ch}^2c^2,
\end{equation*}
质量通过质壳条件进入能量--动量关系。复标量的质量记为 $m_{\rm ch}$，实媒介粒子的静质量记为 $m_\chi$；$m_\chi$ 的 SI 单位为 $\mathrm{kg}$。

由标量场作用量的 SI 归一化有 $[\varphi_{\rm ch}]=[\chi]=\sqrt{\mathrm{J/m}}$。因此三次单项式 $\varphi_{\rm ch}^\dagger\varphi_{\rm ch}\chi$ 的系数不能写成一个无量纲常数。定义具有\emph{动量量纲}的约化三次耦合 $g_3$，并写
\begin{equation}
 \boxed{
 \Lag_{3}
 =-\frac{g_3}{\hbar\sqrt{\hbar c}}\,\varphi_{\rm ch}^\dagger\varphi_{\rm ch}\chi,}
 \qquad [g_3]=\mathrm{kg\,m\,s^{-1}}.
 \label{eq:scalar-cubic-si-r10}
\end{equation}
量纲可逐项检查：$[\hbar\sqrt{\hbar c}]=\mathrm{kg^{3/2}m^{7/2}s^{-2}}$，所以 $[g_3/(\hbar\sqrt{\hbar c})]=\mathrm{J^{-1/2}m^{-3/2}}$；再乘 $[\varphi_{\rm ch}^\dagger\varphi_{\rm ch}\chi]=\mathrm{J^{3/2}m^{-3/2}}$，正好得到 $\mathrm{J/m^3}$。

\begin{exbox}[$\chi\to\varphi_{\rm ch}+\varphi_{\rm ch}^{\ast}$ 的两体衰变]
取母媒介粒子的四动量为 $p_\chi$，末态四动量为 $p_1,p_2$：
\begin{equation*}
 p_\chi=p_1+p_2,\qquad p_\chi^2=m_\chi^2c^2,\qquad p_1^2=p_2^2=m_{\rm ch}^2c^2.
\end{equation*}
只有一个顶角，所以约化振幅为
\begin{equation}
 \ii\mathcal M=-\ii g_3,
 \qquad |\mathcal M|^2=g_3^2.
 \label{eq:scalar-cubic-decay-mphi-r15}
\end{equation}
在母粒子静止系，
\begin{equation*}
 p_\chi^\mu=(m_\chi c,\mathbf0),\qquad
 p_1^\mu=(E/c,\mathbf p),\qquad
 p_2^\mu=(E/c,-\mathbf p).
\end{equation*}
能量守恒给 $2E=m_\chi c^2$。再用 $E^2=m_{\rm ch}^2c^4+c^2p^2$，
\begin{align*}
 p^2
 &=\frac{E^2}{c^2}-m_{\rm ch}^2c^2,\\
 &=\frac{m_\chi^2c^2}{4}-m_{\rm ch}^2c^2,
\end{align*}
所以
\begin{equation}
 \boxed{
 |\mathbf p|
 =\frac{c}{2}\sqrt{m_\chi^2-4m_{\rm ch}^2},
 \qquad m_\chi\ge 2m_{\rm ch}.}
 \label{eq:scalar-cubic-decay-p-r10}
\end{equation}
使用由 Lorentz 不变量二体相空间得到的一体衰变公式 \eqnrefs{eq:master-decay-rate-si-r9}，并代入
\begin{equation*}
 \int\dd\mathfrak P_2
 =\frac{|\mathbf p|}{4\pi m_\chi c},
\end{equation*}
得到
\begin{align}
 \Gamma_{\chi\to\varphi_{\rm ch}\varphi_{\rm ch}^{\ast}}
 &=\frac{1}{2m_\chi\hbar}g_3^2
 \frac{|\mathbf p|}{4\pi m_\chi c},\\
 &=\boxed{
 \frac{g_3^2}{16\pi m_\chi\hbar}
 \sqrt{1-\frac{4m_{\rm ch}^2}{m_\chi^2}}.}
 \label{eq:scalar-cubic-decay-width-r10}
\end{align}
量纲检查为 $[g_3^2/(m_\chi\hbar)]=\mathrm{s^{-1}}$。末态粒子与反粒子可区分，所以没有 $1/2!$；若两个末态是同种实标量，相空间必须再除以 $2!$。


\medskip

\eqnrefs{eq:scalar-cubic-si-r10} 已固定三次相互作用的 SI 归一化。取 Dyson 展开的最低非零阶，把三条外腿接到同一个局域插入点，并按 LSZ 约定截去外腿传播子，就得到\eqnrefs{eq:scalar-cubic-decay-mphi-r15} 中的约化顶角振幅。这里三种场在顶角上彼此可区分，因此不存在额外的 $3!$ 组合因子。

相互作用绘景散射算符的一阶项为
\begin{align*}
\hat S_3^{(1)}
&=\frac{\ii}{\hbar c}\int\dd^4 x\,\Lag_3(x)\\
&=-\frac{\ii g_3}{\hbar^2c\sqrt{\hbar c}}
\int\dd^4 x\,
\varphi_{\rm ch}^\dagger(x)\varphi_{\rm ch}(x)\chi(x).
\end{align*}
把一条 $\varphi_{\rm ch}$、一条 $\varphi_{\rm ch}^\ast$ 和一条 $\chi$ 外腿接到该插入点，树级连通三点函数的极点部分为
\begin{align*}
\widetilde G_{3,c}^{(1)}
={}&(2\pi\hbar)^4\delta^{(4)}(p_\chi-p_1-p_2)
\left(-\frac{\ii g_3}{\hbar^2c\sqrt{\hbar c}}\right)\\
&\times
\widetilde\Delta_{\rm ch}(p_1)
\widetilde\Delta_{\rm ch}(p_2)
\widetilde\Delta_\chi(p_\chi),
\end{align*}
其中
\begin{equation*}
\widetilde\Delta_a(p)
=\frac{\ii\hbar^3c}{p^2-m_a^2c^2+\ii0^+}.
\end{equation*}
树级留数为 $Z_a=1$。对三条外腿分别乘相应逆极点因子并取质量壳极限，三个传播子全部消去，留下
\begin{equation*}
\ii\mathfrak A_3
=-\frac{\ii g_3}{\hbar^2c\sqrt{\hbar c}}.
\end{equation*}
按照约化标量振幅的定义
$\mathcal M_3\equiv\hbar^2c\sqrt{\hbar c}\,\mathfrak A_3$，于是
\begin{equation*}
\ii\mathcal M_3=-\ii g_3,
\end{equation*}
即\eqnrefs{eq:scalar-cubic-decay-mphi-r15} 的第一式。
\end{exbox}

在同一三次耦合理论中，若两种可区分物质场 $\varphi_{{\rm ch},A}$、$\varphi_{{\rm ch},B}$ 分别以 $g_{3A}$、$g_{3B}$ 耦合到同一中性媒介场 $\chi$，最低阶弹性过程
\[
A(p_1)+B(p_2)\to A(p_3)+B(p_4)
\]
由一条 $t$ 道媒介传播线连接两个三点顶角。令交换四动量 $q=p_1-p_3$、$t=q^2$，则树级约化振幅为
\begin{equation}
\boxed{
\mathcal M_t
=\frac{g_{3A}g_{3B}}{m_\chi^2c^2-t}.}
\label{eq:scalar-exchange-mphi-r15}
\end{equation}
分母直接记录媒介传播子的极点位置；$g_{3A}g_{3B}$ 与分母具有相同的动量平方量纲，因此 $\mathcal M_t$ 无量纲。它同时决定相对论散射截面和低速极限中的有效势，两者只是同一交换核在不同运动学区域的读出。

\begin{exbox}[单媒介粒子交换的角分布]
取 $m_A,m_B$ 为两种物质粒子的质量。由正文\eqnrefs{eq:scalar-exchange-mphi-r15}，标量外腿没有自旋求和，故
\begin{equation}
 \boxed{
 \overline{|\mathcal M_t|^2}
 =\frac{g_{3A}^2g_{3B}^2}{(m_\chi^2c^2-t)^2}.}
 \label{eq:scalar-exchange-M2-r10}
\end{equation}
对等质量弹性散射 $m_A=m_B=m_{\rm ch}$，质心系
\begin{equation*}
 s=4(m_{\rm ch}^2c^2+p^2),\qquad
 t=-2p^2(1-\cos\theta),
\end{equation*}
而 $p_f=p_i$。代入 SI 截面主式，
\begin{equation}
 \boxed{
 \frac{\dd\sigma}{\dd\Omega}
 =\frac{\hbar^2}{64\pi^2s}
 \frac{g_{3A}^2g_{3B}^2}
 {[m_\chi^2c^2+2p^2(1-\cos\theta)]^2}.}
 \label{eq:scalar-exchange-angular-r10}
\end{equation}
这里 $\hbar^2/s$ 提供面积量纲，因此该因子必须保留在 SI 截面中。
\end{exbox}

交换振幅还可以回答一个与截面不同的问题：当两粒子运动足够慢、能量转移远小于空间动量转移时，这个相对论交换过程能否被一个定态势能描述？这个问题不能只看传播子分母的形状，因为相对论外态与 Schr\"odinger 两体态采用不同归一化。需要先把同一个低速散射截面在两种描述中匹配，再把动量空间核 Fourier 变换到坐标空间。

令
\begin{equation*}
 m_{\rm red}\equiv\frac{m_A m_B}{m_A+m_B},
 \qquad
 \kappa\equiv\frac{p}{\hbar},
\end{equation*}
分别表示两体相对运动的约化质量与波数。定态 Schr\"odinger 方程可写成 Helmholtz 形式
\begin{equation}
 (\nabla^2+\kappa^2)\psi(\mathbf r)
 =\frac{2m_{\rm red}}{\hbar^2}V(\mathbf r)\psi(\mathbf r),
 \label{eq:sch-helmholtz-r13}
\end{equation}
其出射 Green 核 $G_0^{(+)}$ 给出 Lippmann--Schwinger 方程
\begin{equation}
 \psi(\mathbf r)=\ee^{\ii\boldsymbol\kappa_i\cdot\mathbf r}
 -\frac{2m_{\rm red}}{\hbar^2}\int\dd^3 r'\,
 G_0^{(+)}(\mathbf r-\mathbf r')V(\mathbf r')\psi(\mathbf r').
 \label{eq:sch-lippmann-r13}
\end{equation}
在 $|q^0|/|\mathbf q|\ll1$ 与 $|\mathbf p|/(m_{A,B}c)\ll1$ 的静态低速区间，相对论外态与 Schr\"odinger 两体态的归一化必须先统一。对可区分的弹性二体通道，在前文固定的 $S$ 矩阵、相空间与空间 Fourier 约定下，若 $\mathcal M_{\rm rel}^{\rm NR}(\mathbf q)$ 表示相对论约化振幅的低速矩阵元，则通用匹配关系为
\begin{equation}
\boxed{
\widetilde V_{\rm NR}(\mathbf q)
=-\frac{\hbar^3}{4m_A m_Bc}\,
\mathcal M_{\rm rel}^{\rm NR}(\mathbf q).}
\label{eq:relativistic-to-nr-potential-matching-dp75}
\end{equation}
这里 $\widetilde V_{\rm NR}$ 是对空间坐标作 Fourier 变换得到的势核，量纲为 $\mathrm{J\,m^3}$；若振幅仍含自旋指标，\eqnrefs{eq:relativistic-to-nr-potential-matching-dp75} 按固定自旋通道理解为势算符的矩阵元。这个关系只属于低速、近静态、弹性的一体势描述，不适用于显著能量转移、实粒子产生或不可忽略迟滞的区域。

对当前标量交换，\eqnrefs{eq:scalar-exchange-mphi-r15} 在该极限变成
$\mathcal M_{\rm rel}^{\rm NR}=g_{3A}g_{3B}/(\mathbf q^2+m_\chi^2c^2)$，故\eqnrefs{eq:relativistic-to-nr-potential-matching-dp75} 立即给出
\begin{equation}
 \boxed{
 \widetilde V_{\rm NR}(\mathbf q)
 =-\frac{\hbar^3g_{3A}g_{3B}}{4m_A m_Bc}
 \frac{1}{\mathbf q^2+m_\chi^2c^2}.}
 \label{eq:scalar-yukawa-momentum-kernel-dp68}
\end{equation}
再按全书空间 Fourier 约定变回坐标空间，得到
\begin{equation}
 \boxed{
 V_Y(r)=-\alpha_{AB}^{(s)}\frac{\hbar c}{r}\ee^{-r/\lambda_\chi},
 \qquad
 \alpha_{AB}^{(s)}\equiv\frac{g_{3A}g_{3B}}{16\pi m_A m_Bc^2},
 \qquad
 \lambda_\chi\equiv\frac{\hbar}{m_\chi c}.}
 \label{eq:scalar-yukawa-potential-dp49}
\end{equation}
无量纲 $\alpha_{AB}^{(s)}$ 控制势强度，$\lambda_\chi$ 是媒介粒子的 Compton 长度并控制作用程。$m_\chi\to0$ 时 $\lambda_\chi\to\infty$，屏蔽因子消失；因此有质量传播子的极点位置经三维 Fourier 变换直接生成指数屏蔽与有限作用程。
同一三次顶角允许
\begin{equation*}
 \chi(p_\chi)+\varphi_{\rm ch}(p)\to\chi(p_\chi')+\varphi_{\rm ch}(p').
\end{equation*}
两条外部 $\chi$ 沿物质线有两种次序，因此内部物质线四动量分别为 $p+p_\chi$ 与 $p-p_\chi'$。约化振幅是
\begin{equation}
 \boxed{
 \mathcal M
 =g_3^2\left[
 \frac1{(p+p_\chi)^2-m_{\rm ch}^2c^2+\ii0^+}
 +\frac1{(p-p_\chi')^2-m_{\rm ch}^2c^2+\ii0^+}
 \right].}
 \label{eq:scalar-compton-analogue-si-r10}
\end{equation}
由 $p^2=p'^2=m_{\rm ch}^2c^2$、$p_\chi^2=p_\chi'^2=m_\chi^2c^2$，
\begin{align*}
 (p+p_\chi)^2-m_{\rm ch}^2c^2&=m_\chi^2c^2+2p\cdot p_\chi,\\
 (p-p_\chi')^2-m_{\rm ch}^2c^2&=m_\chi^2c^2-2p\cdot p_\chi'.
\end{align*}
这组公式的目的不是建立一个新的命名模型，而是提前看清后来 Compton 散射会反复出现的拓扑结构：两个外部玻色量子沿同一条物质线连接时，存在两种不同的时间排序或线连接顺序。两项必须先在振幅层相加，再对总振幅取模平方。真正的光子 Compton 散射还会额外包含偏振矢量、Dirac 分子和规范约束；这些结构在 QED 建立后再逐项加入。

树级图的内部动量通常被外部动量守恒完全固定；一旦内部线形成闭合回路，情况就不同了。以 $\varphi^4$ 理论四点函数的 $g_4^2$ 阶 $s$ 道泡形结构为例，两条内部线只要求总四动量为 $P=p_1+p_2$。若第一条内部线取四动量 $\ell$，第二条就是 $P-\ell$，而 $\ell$ 本身仍可连续变化，因此必须对它积分：
\begin{equation}
 \ii\mathcal M_{\rm 1L}^{(s)}
 =
 \frac{(-\ii g_4)^2}{2}
 \int_{\ell}^{(d)}
 \frac{\ii}{\ell^2-m^2c^2+\ii0^+}
 \frac{\ii}{(P-\ell)^2-m^2c^2+\ii0^+}.
 \label{eq:phi4-bubble-toy-r12}
\end{equation}
这种没有被外部动量固定、必须连续积分的内部四动量称为\emph{环动量}，相应闭合内部线结构称为\emph{环图}。这里的 $1/2$ 来自前文 Wick 收缩的组合计数。对 $d=4$，当 $|\ell|\to\infty$ 时，积分核按 $\dd^4 \ell/\ell^4$ 缩放，径向部分表现为 $\int^{\Lambda}\dd\ell/\ell$；因此高动量端出现对数发散。后面所谓正规化与重整化，正是为了先给这类高动量积分一个明确数学定义，再把可测低能参数与调节器依赖分开。




\begin{derivation}[从\eqnrefs{eq:2to2-cross-section-master-si-r9}到\eqnrefs{eq:relativistic-to-nr-potential-matching-dp75}]
低速弹性区有 $|\mathbf p_f|=|\mathbf p_i|$。定义非相对论振幅为阈值极限
\begin{equation*}
\mathcal M_{\rm rel}^{\rm NR}(\mathbf q)
\equiv\lim_{|\mathbf p_i|/(m_{A,B}c)\to0}\mathcal M_{\rm rel},
\qquad
\lim_{|\mathbf p_i|\to0}\sqrt{s}=(m_A+m_B)c.
\end{equation*}
因此对截面主式取同一极限得到
\begin{equation*}
 \frac{\dd\sigma}{\dd\Omega}
 =\frac{\hbar^2}{64\pi^2(m_A+m_B)^2c^2}
 |\mathcal M_{\rm rel}^{\rm NR}(\mathbf q)|^2.
\end{equation*}
另一方面，对\eqnrefs{eq:sch-lippmann-r13} 作一次 Born 迭代，采用同一出射波相位约定可写
\begin{equation*}
 f_B(\mathbf q)
 =-\frac{m_{\rm red}}{2\pi\hbar^2}
 \widetilde V_{\rm NR}(\mathbf q),
 \qquad
 \frac{\dd\sigma}{\dd\Omega}=|f_B|^2,
\end{equation*}
其中 $m_{\rm red}=m_Am_B/(m_A+m_B)$。与相对论截面对应的散射振幅归一化为
\begin{equation*}
 f_{\rm rel}(\mathbf q)
 =\frac{\hbar}{8\pi(m_A+m_B)c}
 \mathcal M_{\rm rel}^{\rm NR}(\mathbf q),
\end{equation*}
其相位由前面固定的 $S$ 矩阵约定与出射 Green 核共同确定。令 $f_B=f_{\rm rel}$，得到
\begin{align*}
 \widetilde V_{\rm NR}(\mathbf q)
 &=-\frac{\hbar^3}{4m_{\rm red}(m_A+m_B)c}
 \mathcal M_{\rm rel}^{\rm NR}(\mathbf q)\\
 &=-\frac{\hbar^3}{4m_Am_Bc}
 \mathcal M_{\rm rel}^{\rm NR}(\mathbf q),
\end{align*}
即\eqnrefs{eq:relativistic-to-nr-potential-matching-dp75}。
\end{derivation}


